\documentclass[UTF8,12pt,a4paper]{ctexart}

\usepackage{geometry}
\geometry{a4paper,left=2.5cm,right=2.5cm,top=3cm,bottom=3cm}

\usepackage{booktabs}
\usepackage{tabularx}
\usepackage{longtable}
\usepackage{array}
\usepackage{xcolor}
\usepackage{graphicx}
\usepackage{amsmath}
\usepackage{amssymb}
\usepackage{listings}
\usepackage{pgfplots}
\pgfplotsset{compat=1.18}
\usepackage{tikz}
\usetikzlibrary{shapes.geometric,arrows.meta,positioning,calc,fit,backgrounds}
\usepackage{tcolorbox}
\tcbuselibrary{skins,breakable}
\usepackage{fancyhdr}
\setlength{\headheight}{14pt}
\usepackage{hyperref}
\usepackage{float}
\usepackage{enumitem}
\usepackage{caption}

% ---- Colors ----
\definecolor{codeblue}{RGB}{0,90,180}
\definecolor{codegray}{RGB}{128,128,128}
\definecolor{codegreen}{RGB}{0,128,64}
\definecolor{codebg}{RGB}{248,248,248}
\definecolor{v9green}{RGB}{40,167,69}
\definecolor{basegray}{RGB}{108,117,125}
\definecolor{warnbg}{RGB}{255,249,230}
\definecolor{warnyellow}{RGB}{220,160,0}

% ---- Listings ----
\lstdefinestyle{cppstyle}{
  language=C++,
  basicstyle=\ttfamily\footnotesize,
  keywordstyle=\color{codeblue}\bfseries,
  commentstyle=\color{codegray}\itshape,
  stringstyle=\color{codegreen},
  numberstyle=\tiny\color{codegray},
  numbers=left,
  numbersep=8pt,
  breaklines=true,
  showstringspaces=false,
  tabsize=4,
  frame=single,
  framerule=0.4pt,
  backgroundcolor=\color{codebg},
  xleftmargin=15pt,
  morekeywords={Kokkos,View,MemoryTraits,Unmanaged,HostSpace,deep_copy,Index_t,Real_t},
}
\lstset{style=cppstyle}

% ---- Header/Footer ----
\pagestyle{fancy}
\fancyhf{}
\fancyhead[L]{\small LULESH 2.0 GPU 移植项目终期报告}
\fancyhead[R]{\small \thepage}
\renewcommand{\headrulewidth}{0.4pt}

% ---- Hyperref ----
\hypersetup{
  colorlinks=true,
  linkcolor=blue!60!black,
  urlcolor=blue!60!black,
  bookmarks=true,
  pdftitle={LULESH 2.0 GPU 移植项目终期报告},
}

\begin{document}

%=============================================================
% Title Page
%=============================================================
\begin{titlepage}
\centering
\vspace*{2.5cm}
{\LARGE\bfseries LULESH 2.0 高性能计算代理应用\par}
\vspace{0.4em}
{\LARGE\bfseries 现代化改造与 GPU 移植\par}
\vspace{0.4em}
{\LARGE\bfseries 终期报告\par}
\vspace{2.5cm}
\begin{tabular}{@{}ll}
  \textbf{框架} & Kokkos 5.0.2（OpenMP / CUDA 后端）、KokkosComm 0.2.0 \\[0.6em]
  \textbf{最终版本} & \texttt{e3e7e98}（2026-04-11，KokkosComm） \\
                    & 集群 CUDA 构建：\texttt{f2520fa} \\[0.6em]
  \textbf{报告日期} & 2026-04-30 \\
\end{tabular}
\vfill
\end{titlepage}

\tableofcontents
\newpage

%=============================================================
\section{移植过程}
%=============================================================

\subsection{项目背景}

LULESH 2.0（Livermore Unstructured Lagrangian Explicit Shock Hydrodynamics）是 LLNL 发布的 HPC
代理应用，其原始代码为单文件（约 2813 行）、MPI + \texttt{\#pragma omp} 混合实现，不支持 GPU。
本项目目标是以 Kokkos~5 为统一并行抽象，完成从 CPU-only OpenMP 到 CPU+GPU 可移植实现的全流程改造，
并在集群节点上验证最终版本相对官方基线的性能提升。

\textbf{版本序列如表~\ref{tab:versions} 所示。}

\begin{table}[H]
\centering
\caption{项目版本序列}
\label{tab:versions}
\begin{tabular}{llll}
\toprule
\textbf{版本} & \textbf{提交哈希} & \textbf{日期} & \textbf{说明} \\
\midrule
baseline & \texttt{3e01c40} & ---        & 官方 LLNL LULESH 2.0.3 \\
v1       & \texttt{3003bb6} & 2026-03-03 & 模块化重构 + Kokkos OpenMP \\
v2       & \texttt{0b4c125} & 2026-03-04 & Kokkos OpenMP 初版 \\
v3       & \texttt{4d14004} & 2026-03-09 & Kokkos OpenMP 优化版 \\
\textbf{v9} & \textbf{\texttt{e3e7e98}} & \textbf{2026-04-11} & \textbf{KokkosComm + P1--P6（最终版）} \\
\bottomrule
\end{tabular}
\end{table}

\subsection{移植路径总览}

\begin{figure}[H]
\centering
\begin{tikzpicture}[
  node distance=1.8cm,
  block/.style={
    rectangle, rounded corners=4pt, draw=blue!50, fill=blue!7,
    text width=8.5cm, align=center, inner sep=8pt, font=\small
  },
  startblock/.style={
    rectangle, rounded corners=4pt, draw=gray!60, fill=gray!10,
    text width=8.5cm, align=center, inner sep=8pt, font=\small
  },
  endblock/.style={
    rectangle, rounded corners=4pt, draw=green!60!black, fill=green!10,
    text width=8.5cm, align=center, inner sep=8pt, font=\small, line width=1.2pt
  },
  arr/.style={-{Stealth[length=7pt]}, thick, draw=blue!60},
  lbl/.style={text width=5.5cm, align=left, font=\scriptsize, text=gray!70!black}
]

\node[startblock] (A) {
  \textbf{原始 LULESH 2.0 (baseline)}\\
  单文件 \texttt{lulesh.cc} $\cdot$
  \texttt{\#pragma omp parallel for}\\
  \texttt{std::vector} / 裸指针 $\cdot$ MPI 混合并行
};
\node[block, below=2.4cm of A] (B) {
  \textbf{v1--v3 \quad Kokkos OpenMP}\\
  10 个物理模块 $\cdot$ \texttt{Kokkos::View} 统一内存\\
  内核融合 + region 展平
};
\node[block, below=2.4cm of B] (C) {
  \textbf{Kokkos GPU 就绪版}\\
  CPU + CUDA 双路径构建 $\cdot$
  \texttt{KOKKOS\_LAMBDA} 核函数
};
\node[endblock, below=2.4cm of C] (D) {
  \textbf{v9 \quad \texttt{e3e7e98} / \texttt{f2520fa}}\\
  KokkosComm 三层通信 $\cdot$ CPU + GPU 双路径\\
  P1--P6 预分配优化
};

\draw[arr] (A) -- (B);
\draw[arr] (B) -- (C);
\draw[arr] (C) -- (D);

\node[lbl, right=0.6cm of A, yshift=-1.2cm] {
  W10: 模块化 + C++20\\
  \texttt{parallel\_for} 替换 \texttt{omp pragma}\\
  两轮内核融合（减少 $\sim$27 barriers/步）
};
\node[lbl, right=0.6cm of B, yshift=-1.2cm] {
  W12: \texttt{KOKKOS\_INLINE\_FUNCTION}\\
  \texttt{View} 化 \texttt{EosTemps}/\texttt{vnew}/\texttt{determ}\\
  W13: \texttt{KOKKOS\_LAMBDA} 迁移（20+ 核函数）
};
\node[lbl, right=0.6cm of C, yshift=-1.2cm] {
  W14: 3D 笛卡尔域分解 / 6 面 halo 交换\\
  W15: KokkosComm 三层架构 + comm dup\\
  W16: P1--P6 \texttt{malloc} 消除
};

\end{tikzpicture}
\caption{LULESH 2.0 移植路径演进}
\label{fig:evolution}
\end{figure}

\subsection{主要技术改造}

\subsubsection{内核融合（W10--W11）}

原始 Kokkos 迁移将每个 \texttt{omp parallel for} 替换为独立 \texttt{parallel\_for}，
但 Kokkos 每个 \texttt{parallel\_for} 末尾有隐式全量 barrier，而原始 LULESH 大量使用
\texttt{omp nowait} 绕过 barrier。通过 macOS \texttt{sample} 工具分析，
barrier 开销占运行时间 \textbf{79.5\%}。

两轮共 11 项内核融合，将 barrier 数量从每步约 50+ 降至约 23（表~\ref{tab:fusion}）。

\begin{table}[H]
\centering
\caption{内核融合轮次与收益}
\label{tab:fusion}
\begin{tabular}{lll}
\toprule
\textbf{轮次} & \textbf{典型融合} & \textbf{节省 barrier} \\
\midrule
第一轮（3 项） & EvalEOS 准备 5$\to$1、CalcEnergy 11$\to$1 & $\sim$15/region/rep \\
第二轮（8 项） & 双 reduce、scatter+声速、速度+位置 & $\sim$20/step \\
\bottomrule
\end{tabular}
\end{table}

同步优化：EOS 临时向量预分配（消除每步约 15 万次 \texttt{malloc}）、
scatter buffer 64 字节对齐（消除 false sharing）、
黏性/时间步 region 循环展平（11$\to$1 个 kernel）。

\subsubsection{GPU 基础设施（W12--W13）}

关键改动：
\begin{itemize}[noitemsep,topsep=2pt]
  \item 几何辅助函数添加 \texttt{KOKKOS\_INLINE\_FUNCTION}
        （即 \texttt{\_\_host\_\_ \_\_device\_\_}）
  \item \texttt{EosTemps}、\texttt{ScatterBuffers}、\texttt{vnew}/\texttt{determ}
        从 \texttt{std::vector} 迁移至 \texttt{Kokkos::View<Real\_t*>}
  \item 全部 \texttt{parallel\_for} 的 \texttt{[\&]} 捕获改为 \texttt{KOKKOS\_LAMBDA}
  \item GPU 兼容性修复：\texttt{EvalEOSForElems} 中区域元素索引列表使用编译期双路径，
        CPU 路径零拷贝，GPU 路径显式 H$\to$D 传输
\end{itemize}

\begin{lstlisting}[caption={EvalEOSForElems 双路径 GPU 修复（src/lulesh-eos.cc）},
                   label={lst:gpufix}]
#ifdef KOKKOS_ENABLE_CUDA
    // GPU: regElemRaw lives in HostSpace; allocate device memory and copy
    Kokkos::View<const Index_t*, Kokkos::HostSpace,
                 Kokkos::MemoryTraits<Kokkos::Unmanaged>>
        regElemList_h(regElemRaw, numElemReg);
    Kokkos::View<Index_t*> regElemList("regElemList", numElemReg);
    Kokkos::deep_copy(regElemList, regElemList_h);
#else
    // CPU: zero-copy, wrap host pointer directly
    Kokkos::View<Index_t*, Kokkos::HostSpace,
                 Kokkos::MemoryTraits<Kokkos::Unmanaged>>
        regElemList(regElemRaw, numElemReg);
#endif
\end{lstlisting}

\subsubsection{KokkosComm 三层通信架构（W14--W15）}

\begin{figure}[H]
\centering
\begin{tikzpicture}[
  layer/.style={rectangle, rounded corners=3pt, draw, minimum width=3.2cm,
                minimum height=3.2cm, align=center, inner sep=6pt, font=\small},
  arr/.style={-{Stealth[length=6pt]}, thick},
]
\node[layer, fill=gray!10, draw=gray!60] (C) at (0,0) {
  \textbf{C 层}\\（业务代码）\\[4pt]
  \texttt{lulesh-comm.cc}\\
  \texttt{lulesh-timestep.cc}\\
  \texttt{lulesh-nodal.cc}\\
  等所有业务文件
};
\node[layer, fill=blue!7, draw=blue!50] (B) at (4.5,0) {
  \textbf{B 层}\\（运行时接口）\\[4pt]
  \texttt{lulesh-runtime.h}\\
  \texttt{WorldComm()}\\
  \texttt{DistributedRank()}\\
  \texttt{PostDtIallreduce()}
};
\node[layer, fill=blue!14, draw=blue!60] (A) at (9,0) {
  \textbf{A 层}\\（MPI 适配）\\[4pt]
  \texttt{lulesh-runtime-mpi.cc}\\
  唯一含 \texttt{\#include mpi.h}\\
  \texttt{MPI\_Init}\\
  \texttt{MPI\_Comm\_dup}
};
\node[layer, fill=orange!10, draw=orange!60, minimum height=2.2cm] (KC) at (4.5,-3.6) {
  \textbf{KokkosComm}\\[4pt]
  \texttt{send}/\texttt{recv}\\
  \texttt{allreduce}\\
  \texttt{Communicator}
};

\draw[arr] (C.east) -- node[above,font=\scriptsize]{调用抽象接口} (B.west);
\draw[arr] (B.east) -- node[above,font=\scriptsize]{实现} (A.west);
\draw[arr] (C.south) to[bend right=15] node[left,font=\scriptsize]{直接调用} (KC.west);
\draw[arr] (KC.east) to[bend right=15] node[right,font=\scriptsize]{底层路由} (A.south);
\end{tikzpicture}
\caption{KokkosComm 三层通信架构}
\label{fig:kokkoscomm}
\end{figure}

\textbf{KokkosComm tag 问题}：KokkosComm 0.2.0 点对点接口固定使用 tag=17，
无法区分两类 halo 通信相位（SBN / MonoQ）。
解决方案：为每个相位单独 \texttt{duplicate} 一个 MPI communicator——
MPI 消息匹配同时检查 \texttt{(source, tag, communicator)}，不同 comm 的消息绝对不会串相。

每步通信数据流如图~\ref{fig:commflow} 所示。

\begin{figure}[H]
\centering
\begin{tikzpicture}[
  ptcp/.style={rectangle, draw=blue!50, fill=blue!8,
               minimum width=2.8cm, minimum height=0.65cm,
               align=center, font=\small\bfseries},
  fwd/.style={-{Stealth[length=5pt]}, thick, draw=blue!70},
  ret/.style={-{Stealth[length=5pt]}, thick, draw=green!60!black, dashed},
  lbl/.style={font=\scriptsize, above},
  note/.style={rectangle, draw=gray!40, fill=yellow!8,
               font=\scriptsize, inner sep=3pt, align=center},
  life/.style={dashed, gray!45, line width=0.8pt},
]
% x positions
\def\xG{0} \def\xH{5} \def\xM{10}
% participant headers
\node[ptcp] (pG) at (\xG,0)  {GPU 计算};
\node[ptcp] (pH) at (\xH,0)  {Host 缓冲区};
\node[ptcp] (pM) at (\xM,0)  {MPI 通信层};
% lifelines
\draw[life] (\xG,-0.35) -- (\xG,-11.5);
\draw[life] (\xH,-0.35) -- (\xH,-11.5);
\draw[life] (\xM,-0.35) -- (\xM,-11.5);
% messages
\draw[fwd] (\xG,-1) -- (\xH,-1) node[lbl,midway]{\texttt{deep\_copy(fx/fy/fz)}};
\draw[fwd] (\xH,-2) -- (\xM,-2) node[lbl,midway]{\texttt{CommRecv} (Irecv 预投递)};
\draw[fwd] (\xH,-3) -- (\xM,-3) node[lbl,midway]{\texttt{CommSend} (Isend 非阻塞)};
\node[note] at (7.5,-3.9) {通信飞行中};
\draw[ret] (\xM,-4.8) -- (\xH,-4.8) node[lbl,midway]{\texttt{CommSBN} (Wait + scatter-add)};
\draw[fwd] (\xH,-5.8) -- (\xG,-5.8) node[lbl,midway]{\texttt{deep\_copy(fx/fy/fz 更新后)}};
\draw[fwd] (\xG,-6.8) -- (\xH,-6.8) node[lbl,midway]{\texttt{deep\_copy(delv\_xi/eta/zeta)}};
\draw[fwd] (\xH,-7.8) -- (\xM,-7.8) node[lbl,midway]{\texttt{CommMonoQ} (发送速度梯度)};
\draw[ret] (\xM,-8.8) -- (\xH,-8.8) node[lbl,midway]{Wait + 写入 ghost 槽};
\draw[fwd] (\xH,-9.8) -- (\xG,-9.8) node[lbl,midway]{\texttt{deep\_copy(ghost 数据)}};
\node[note, text width=7cm] at (5,-11)
      {\texttt{MPI\_Allreduce(MIN)} 全局时间步归约};
\end{tikzpicture}
\caption{每步 GPU--MPI 通信数据流}
\label{fig:commflow}
\end{figure}

\subsubsection{代码级优化 P1--P6（W16）}

\begin{table}[H]
\centering
\caption{P1--P6 代码级优化}
\label{tab:p1p6}
\begin{tabular}{lll}
\toprule
\textbf{编号} & \textbf{改动} & \textbf{消除开销} \\
\midrule
P1 & 删除冗余体积检查内核 & 1 次 kernel launch/步 \\
P3 & 预分配 SBN host mirrors & $\sim$715 KB malloc/步 \\
P4 & 预分配 vnew + determ & $\sim$432 KB malloc/步 \\
\textbf{P5} & \textbf{预分配 scatter 缓冲区（6 个 View）} & \textbf{$\sim$10.4 MB malloc/步（最大项）} \\
P6 & 融合 ApplyBC 三内核 & 2 次 launch/步 \\
\bottomrule
\end{tabular}
\end{table}

%=============================================================
\section{性能测试与分析}
%=============================================================

\subsection{测试环境}

全部基准测试在集群节点 \texttt{cpu-1239} / \texttt{gpu-1239} 上完成，保证版本间对比公平。

\begin{table}[H]
\centering
\caption{测试环境概况}
\label{tab:testenv}
\begin{tabular}{ll}
\toprule
\textbf{项目} & \textbf{说明} \\
\midrule
集群节点 & \texttt{cpu-1239}（CPU 测试）；\texttt{gpu-1239}（GPU 测试） \\
线程配置 & T=1 / 2 / 4 / 8 \\
每组重复 & 5 次，取中位数 \\
性能指标 & \textbf{FOM}（zone-updates/s）= 总单元更新次数 / 运行时间 \\
参考基线 & \texttt{cpu-baseline/3e01c40}（官方 LLNL LULESH 2.0.3，同节点构建运行） \\
\bottomrule
\end{tabular}
\end{table}

\begin{table}[H]
\centering
\caption{测试规模配置}
\label{tab:testscale}
\begin{tabular}{lll}
\toprule
\textbf{类型} & \textbf{参数} & \textbf{说明} \\
\midrule
强扩展（小） & \texttt{-s 30 -i 100} & 27,000 单元，问题规模固定，线程倍增 \\
强扩展（大） & \texttt{-s 45 -i 200} & 91,125 单元 \\
弱扩展 & T=1:s30, T=2:s38, T=4:s48, T=8:s60 & 每线程工作量$\approx$恒定（zones $\propto$ T） \\
\bottomrule
\end{tabular}
\end{table}

\subsection{正确性验证}

所有版本通过正确性基准 \texttt{-s 10 -i 50}
（参考值 \texttt{Final Origin Energy = 8.104796e+04}）：

\begin{table}[H]
\centering
\caption{正确性验证结果}
\label{tab:correctness}
\begin{tabular}{lllll}
\toprule
\textbf{版本} & \textbf{配置} & \textbf{Final Energy} &
\textbf{MaxRelDiff} & \textbf{结果} \\
\midrule
baseline   & T=1    & $8.104796\times10^{4}$ & $2.92\times10^{-13}$ & $\checkmark$ \\
v9 CPU     & T=1    & $8.104796\times10^{4}$ & $2.92\times10^{-13}$ & $\checkmark$ \\
v9 GPU     & 1 GPU  & $8.104796\times10^{4}$ & $1.30\times10^{-13}$ & $\checkmark$ \\
v9 MPI 8进程 & n=8  & $6.483837\times10^{5}$ & $<1\times10^{-10}$   & $\checkmark$ \\
\bottomrule
\end{tabular}
\end{table}

此外，v9 完成了 GitHub 官方 LLNL LULESH 2.0.3 基线（commit \texttt{46c2a1d6}）的
\textbf{C1--C11 共 22 次双边运行}全面对照，所有 \texttt{Final Origin Energy}
与迭代次数完全一致。

\subsection{版本演进性能（CPU 强扩展，\texttt{-s 30 -i 100}）}

\begin{table}[H]
\centering
\caption{强扩展 FOM 绝对值（中位数，zones/s，\texttt{-s 30 -i 100}）}
\label{tab:strong30}
\begin{tabular}{lrrrr}
\toprule
\textbf{版本} & \textbf{T=1} & \textbf{T=2} & \textbf{T=4} & \textbf{T=8} \\
\midrule
baseline (\texttt{3e01c40}) & 1,519 & 2,174 & 3,499 & 4,621 \\
v1 重构 OpenMP              & 1,363 & 1,616 & 2,012 & 2,109 \\
v2 Kokkos 初版              & 1,146 & 1,448 & 1,868 & 2,035 \\
v3 Kokkos 优化              & 1,655 & 2,293 & 4,011 & 6,135 \\
\textbf{v9 KokkosComm}      & \textbf{1,742} & \textbf{3,175} & \textbf{5,857} & \textbf{9,831} \\
\bottomrule
\end{tabular}
\end{table}

\begin{figure}[H]
\centering
\begin{tikzpicture}
\begin{axis}[
  ybar, width=12cm, height=7cm, bar width=22pt,
  symbolic x coords={baseline, v1, v2, v3, v9},
  xtick=data,
  xticklabel style={font=\small},
  ylabel={FOM (zones/s)}, ylabel style={font=\small},
  ymin=0, ymax=11500,
  ytick={0,2000,4000,6000,8000,10000},
  ymajorgrids=true, grid style={dashed,gray!30},
  nodes near coords, nodes near coords style={font=\scriptsize,rotate=90,anchor=west},
  title={各版本 T=8 FOM（\texttt{-s 30 -i 100}，中位数）},
  title style={font=\small},
]
\addplot[fill=basegray!70, draw=none] coordinates
  {(baseline,4621)(v1,2109)(v2,2035)(v3,6135)(v9,9831)};
\end{axis}
\end{tikzpicture}
\caption{各版本 T=8 FOM 演进（\texttt{-s 30 -i 100}）}
\label{fig:fom30}
\end{figure}

\begin{table}[H]
\centering
\caption{强扩展效率（$\mathrm{FOM}_T / (T \times \mathrm{FOM}_{T=1})$，\texttt{-s 30 -i 100}）}
\label{tab:eff30}
\begin{tabular}{lrrr}
\toprule
\textbf{版本} & \textbf{T=2} & \textbf{T=4} & \textbf{T=8} \\
\midrule
baseline     & 71.5\% & 57.6\% & 38.0\% \\
v1           & 59.3\% & 36.9\% & 19.3\% \\
v2           & 63.1\% & 40.7\% & 22.2\% \\
v3           & 69.3\% & 60.6\% & 46.3\% \\
\textbf{v9}  & \textbf{91.1\%} & \textbf{84.1\%} & \textbf{70.6\%} \\
\bottomrule
\end{tabular}
\end{table}

\begin{figure}[H]
\centering
\begin{tikzpicture}
\begin{axis}[
  xbar, width=11cm, height=5.5cm, bar width=14pt,
  symbolic y coords={v9,v3,v2,v1,baseline},
  ytick=data,
  yticklabel style={font=\small},
  xlabel={强扩展效率 (\%)}, xlabel style={font=\small},
  xmin=0, xmax=100,
  xtick={0,20,40,60,80,100},
  xmajorgrids=true, grid style={dashed,gray!30},
  nodes near coords, nodes near coords style={font=\scriptsize},
  nodes near coords align={horizontal},
  title={T=8 强扩展效率对比（\texttt{-s 30 -i 100}）},
  title style={font=\small},
]
\addplot[fill=blue!35, draw=none] coordinates
  {(38.0,baseline)(19.3,v1)(22.2,v2)(46.3,v3)(70.6,v9)};
\end{axis}
\end{tikzpicture}
\caption{T=8 强扩展效率对比（\texttt{-s 30 -i 100}）}
\label{fig:eff30bar}
\end{figure}

\textbf{关键观察：}
\begin{itemize}[noitemsep,topsep=2pt]
  \item \textbf{v1/v2 效率低于 baseline}（T=8：19--22\% vs 38\%）：
        原始 LULESH 使用 \texttt{omp nowait} 允许快核心不等待慢核心；
        Kokkos \texttt{parallel\_for} 的隐式全量 barrier 消除了这一优化，
        使 Apple Silicon P/E 异构核心问题完全暴露。
  \item \textbf{v3 超越 baseline}（T=8 FOM: 6,135 vs 4,621，+32.8\%）：
        两轮内核融合使 T=4+ 效率超越 baseline。
  \item \textbf{v9 大幅领先}（T=8 FOM: 9,831 vs 4,621，\textbf{+112.7\%}）：
        P1--P6 \texttt{malloc} 消除 + 集群均质 CPU 核心（无 P/E 干扰）+
        KokkosComm 通信层共同作用。
\end{itemize}

\subsection{版本演进性能（CPU 强扩展，\texttt{-s 45 -i 200}）}

\begin{table}[H]
\centering
\caption{强扩展 FOM（中位数，zones/s，\texttt{-s 45 -i 200}）}
\label{tab:strong45}
\begin{tabular}{lrrrrr}
\toprule
\textbf{版本} & \textbf{T=1} & \textbf{T=2} & \textbf{T=4} & \textbf{T=8} & \textbf{T=8 效率} \\
\midrule
baseline    & 1,594 & 2,249 & 3,441 & 4,222 & 33.1\% \\
v1          & 1,410 & 1,595 & 1,885 & 1,965 & 17.4\% \\
v3          & 1,493 & 2,379 & 3,904 & 5,454 & 45.8\% \\
\textbf{v9} & \textbf{1,858} & \textbf{3,507} & \textbf{5,836} & \textbf{7,813} & \textbf{52.5\%} \\
\bottomrule
\end{tabular}
\end{table}

v9 T=8 FOM 7,813 vs baseline 4,222，提升 \textbf{+85.1\%}；
T=8 效率 52.5\%（较 \texttt{-s 30} 的 70.6\% 略降，与大问题下通信比例上升一致）。

\subsection{弱扩展性能}

\begin{table}[H]
\centering
\caption{弱扩展 FOM（中位数，zones/s）}
\label{tab:weak}
\begin{tabular}{lrrrr}
\toprule
\textbf{版本} & \textbf{T=1 (s30)} & \textbf{T=2 (s38)} & \textbf{T=4 (s48)} & \textbf{T=8 (s60)} \\
\midrule
baseline    & 1,479 & 2,220 & 3,536 & 4,102 \\
v1          & 1,329 & 1,604 & 1,880 & 1,923 \\
v3          & 1,608 & 2,389 & 4,030 & 5,197 \\
\textbf{v9} & \textbf{1,692} & \textbf{3,403} & \textbf{5,926} & \textbf{7,574} \\
\bottomrule
\end{tabular}
\end{table}

\begin{table}[H]
\centering
\caption{弱扩展效率（$\mathrm{FOM}_T / (T \times \mathrm{FOM}_{T=1})$，理想值 = 100\%）}
\label{tab:weakeff}
\begin{tabular}{lrrr}
\toprule
\textbf{版本} & \textbf{T=2} & \textbf{T=4} & \textbf{T=8} \\
\midrule
baseline     & 75.1\% & 59.8\% & 34.7\% \\
v1           & 60.4\% & 35.4\% & 18.1\% \\
v3           & 74.3\% & 62.7\% & 40.4\% \\
\textbf{v9}  & \textbf{100.6\%} & \textbf{87.6\%} & \textbf{55.9\%} \\
\bottomrule
\end{tabular}
\end{table}

\begin{figure}[H]
\centering
\begin{tikzpicture}
\begin{axis}[
  xbar, width=11cm, height=5cm, bar width=14pt,
  symbolic y coords={v9,v3,v1,baseline},
  ytick=data,
  yticklabel style={font=\small},
  xlabel={弱扩展效率 (\%)}, xlabel style={font=\small},
  xmin=0, xmax=75,
  xtick={0,20,40,60},
  xmajorgrids=true, grid style={dashed,gray!30},
  nodes near coords, nodes near coords style={font=\scriptsize},
  nodes near coords align={horizontal},
  title={T=8 弱扩展效率对比},
  title style={font=\small},
]
\addplot[fill=v9green!50, draw=none] coordinates
  {(34.7,baseline)(18.1,v1)(40.4,v3)(55.9,v9)};
\end{axis}
\end{tikzpicture}
\caption{T=8 弱扩展效率对比}
\label{fig:weakeff}
\end{figure}

v9 T=2 弱扩展效率 100.6\%，接近理想线性扩展；
T=8 效率 55.9\%，领先 baseline（34.7\%）和 v3（40.4\%）。

\subsection{v9 与官方 baseline 性能对比汇总}

\begin{table}[H]
\centering
\caption{v9 vs baseline 性能对比（集群节点 \texttt{cpu-1239}）}
\label{tab:v9baseline}
\begin{tabular}{lrrr}
\toprule
\textbf{测试配置} & \textbf{baseline FOM} & \textbf{v9 FOM} & \textbf{加速比} \\
\midrule
T=1, \texttt{-s 30 -i 100}          & 1,519 & 1,742 & 1.15$\times$ \\
T=2, \texttt{-s 30 -i 100}          & 2,174 & 3,175 & 1.46$\times$ \\
T=4, \texttt{-s 30 -i 100}          & 3,499 & 5,857 & 1.67$\times$ \\
\textbf{T=8, \texttt{-s 30 -i 100}} & \textbf{4,621} & \textbf{9,831} & \textbf{2.13$\times$} \\
T=1, \texttt{-s 45 -i 200}          & 1,594 & 1,858 & 1.17$\times$ \\
T=4, \texttt{-s 45 -i 200}          & 3,441 & 5,836 & 1.70$\times$ \\
\textbf{T=8, \texttt{-s 45 -i 200}} & \textbf{4,222} & \textbf{7,813} & \textbf{1.85$\times$} \\
\bottomrule
\end{tabular}
\end{table}

\begin{figure}[H]
\centering
\begin{tikzpicture}
\begin{axis}[
  ybar, width=13cm, height=7cm, bar width=22pt,
  symbolic x coords={T1 s30, T4 s30, T8 s30, T1 s45, T4 s45, T8 s45},
  xtick=data,
  xticklabel style={font=\small},
  ylabel={FOM 提升幅度 (\%)}, ylabel style={font=\small},
  ymin=0, ymax=130,
  ytick={0,20,40,60,80,100,120},
  ymajorgrids=true, grid style={dashed,gray!30},
  nodes near coords, nodes near coords style={font=\scriptsize},
  title={v9 相对 baseline 的 FOM 提升幅度},
  title style={font=\small},
]
\addplot[fill=v9green!65, draw=none] coordinates
  {(T1 s30,14.7)(T4 s30,67.4)(T8 s30,112.7)
   (T1 s45,16.6)(T4 s45,69.6)(T8 s45,85.1)};
\end{axis}
\end{tikzpicture}
\caption{v9 相对 baseline 的 FOM 提升幅度（\%）}
\label{fig:improvement}
\end{figure}

\textbf{关键结论}：v9 在 T=8 时较 baseline 快约 \textbf{2$\times$}。
加速比随线程数增加而扩大（T=1 时约 1.15$\times$，T=8 时约 2.13$\times$），
说明性能提升主要来源于更优的并行扩展性（T=8 效率 70.6\% vs baseline 38.0\%），
而非单线程的算法改进。

\subsection{GPU 测试（简述）}

v9 的 CUDA 构建（\texttt{f2520fa}）在集群节点 \texttt{gpu-1239} 单 GPU 上运行的基准数据
如表~\ref{tab:gpu} 所示。

\begin{table}[H]
\centering
\caption{GPU 性能数据（单 GPU，\texttt{gpu-1239}）}
\label{tab:gpu}
\begin{tabular}{lrrr}
\toprule
\textbf{测试规模} & \textbf{GPU FOM} & \textbf{vs v9 CPU T=8} & \textbf{vs baseline T=8} \\
\midrule
\texttt{-s 30 -i 100} & 14,355 & 1.46$\times$ & 3.11$\times$ \\
\texttt{-s 45 -i 200} & 23,327 & 2.99$\times$ & 5.52$\times$ \\
\bottomrule
\end{tabular}
\end{table}

单 GPU 性能约为 8 线程 CPU v9 的 1.5--3.0$\times$，约为 baseline T=8 的 3--5.5$\times$。
GPU 在大问题规模（\texttt{-s 45}）下优势更显著，符合 GPU 的大规模并行特性。

\begin{tcolorbox}[
  colback=gray!5, colframe=gray!45,
  fontupper=\small, left=4pt, right=4pt, top=3pt, bottom=3pt]
  \textbf{注：}当前 GPU 测试为单 GPU 单进程，多 GPU MPI 协同测试尚未完成。
\end{tcolorbox}

%=============================================================
\section{集群完整性能对比（待补充）}
%=============================================================

\begin{tcolorbox}[
  colback=warnbg, colframe=warnyellow,
  fontupper=\small, left=6pt, right=6pt, top=5pt, bottom=5pt,
  title={\textbf{本节结果尚未可用}}]
集群上最终版本（\texttt{e3e7e98} / \texttt{f2520fa}）与官方 LULESH 2.0.3
的完整对比测试正在进行中，测试结果出来后将补充本节内容。
\end{tcolorbox}

\bigskip
本节计划涵盖以下对比内容：

\begin{table}[H]
\centering
\caption{计划测试矩阵}
\label{tab:planned}
\begin{tabular}{lll}
\toprule
\textbf{测试类型} & \textbf{参数范围} & \textbf{说明} \\
\midrule
CPU 单节点强扩展 & \texttt{-s 30,45} / T=1,2,4,8 & v9 CPU vs baseline，主要对比指标 \\
CPU 弱扩展       & T=1--8，等比放大规模           & 评估扩展效率 \\
MPI 多进程       & n=1,8,27                       & 验证分布式并行正确性与效率 \\
GPU vs CPU baseline & 1 GPU vs baseline T=8       & 量化 GPU 加速效益 \\
\bottomrule
\end{tabular}
\end{table}

\textbf{预期分析维度：}
\begin{itemize}[noitemsep,topsep=2pt]
  \item 强扩展效率曲线（Amdahl 定律验证）
  \item 各优化阶段（内核融合 / \texttt{malloc} 消除 / KokkosComm）的独立贡献拆分
  \item 通信开销占比（基于 profiling：SBN wait、MonoQ wait、dt allreduce 各自份额）
\end{itemize}

%=============================================================
\section{总结与展望}
%=============================================================

\subsection{工作成果}

本项目历经 7 周完成了 LULESH 2.0 从原始 OpenMP 代码到
Kokkos~5 + KokkosComm GPU-ready 实现的全流程改造，主要成果：

\textbf{可移植性}：同一代码库支持 OpenMP 和 CUDA 两种后端，无需维护两套实现。
\texttt{KOKKOS\_LAMBDA} 和 \texttt{\#ifdef KOKKOS\_ENABLE\_CUDA} 编译期分支
兼顾 CPU 零拷贝性能与 GPU 正确性。

\textbf{正确性}：通过官方 LLNL LULESH 2.0.3 基线的 C1--C11 共 22 次双边验证，
以及 1/8/27 进程等价性测试，GPU 版本数值精度与 CPU 版本一致
（MaxRelDiff $< 10^{-13}$）。

\textbf{CPU 性能}（集群节点，与官方 baseline 对比）：

\begin{table}[H]
\centering
\caption{核心性能指标对比汇总}
\label{tab:summary}
\begin{tabular}{llll}
\toprule
\textbf{指标} & \textbf{baseline} & \textbf{v9} & \textbf{提升} \\
\midrule
T=8 FOM（\texttt{-s 30}）          & 4,621   & 9,831  & \textbf{+112.7\%} \\
T=8 强扩展效率（\texttt{-s 30}）   & 38.0\%  & 70.6\% & \textbf{+32.6 pp} \\
T=8 弱扩展效率（\texttt{-s 60}）   & 34.7\%  & 55.9\% & \textbf{+21.2 pp} \\
\bottomrule
\end{tabular}
\end{table}

\textbf{工程质量}：代码从 2813 行单文件重构为 10 个物理模块；
KokkosComm 三层通信架构完全隔离 MPI 依赖；
修复 5 项历史缺陷（MPI\_Abort 安全性、GPU kernel 死代码等）。

\subsection{局限性}

\textbf{多 GPU MPI 协同}：通信点当前使用显式 \texttt{deep\_copy}
（GPU$\to$host$\to$MPI$\to$host$\to$GPU），CUDA-aware MPI
可直接传递 device 指针消除中间拷贝，是多节点 GPU 运行效率的关键缺口。

\textbf{通信计算重叠（T1/T2 优化）}：W16 的实验已将 plane recv 暴露比从 0.65 降至
0.03，但 release A/B 测试尚未完成，wall-time 收益待量化。

\textbf{EOS region 索引 H$\to$D 拷贝}：每步对 GPU 路径执行 \texttt{numReg}
次小规模 H$\to$D 传输，可通过将 \texttt{m\_regElemlist}
迁移为设备端平铺 View 消除。

\textbf{KokkosComm 0.2.0 局限}：\texttt{wait\_all} 不会重置 request 对象，
需逐个 \texttt{wait}；顶层 P2P 接口不暴露 tag 参数，
通过 communicator duplication 绕过但增加了 MPI 资源开销。

\subsection{未来工作}

\begin{enumerate}[noitemsep,topsep=2pt]
  \item \textbf{CUDA-aware MPI}：直接传递 device 指针，
        消除 GPU$\leftrightarrow$host 双向拷贝，是多 GPU 性能的首要瓶颈。
  \item \textbf{T2 正式 A/B 测量}：完成通信重叠的 release 性能验证，
        若效果显著推广至 SBN 阶段。
  \item \textbf{多 GPU 扩展性测试}：在 GPU 节点以 n=1,2,4,8 GPU
        测试强/弱扩展，验证 GPU 路径多节点正确性。
  \item \textbf{EOS 索引 View 化}：消除 per-step H$\to$D 拷贝，
        对多 region 场景（\texttt{-r 11}）GPU 路径有显著收益。
\end{enumerate}

%=============================================================
\appendix
\section{完整性能数据表}
%=============================================================

\subsection{强扩展中位 FOM（\texttt{-s 30 -i 100}，集群 \texttt{cpu-1239}）}

\begin{table}[H]
\centering
\begin{tabular}{lrrrr}
\toprule
\textbf{版本} & \textbf{T=1} & \textbf{T=2} & \textbf{T=4} & \textbf{T=8} \\
\midrule
baseline (\texttt{3e01c40}) & 1,519 & 2,174 & 3,499 & 4,621 \\
v1 (\texttt{3003bb6})       & 1,363 & 1,616 & 2,012 & 2,109 \\
v2 (\texttt{0b4c125})       & 1,146 & 1,448 & 1,868 & 2,035 \\
v3 (\texttt{4d14004})       & 1,655 & 2,293 & 4,011 & 6,135 \\
v9 (\texttt{e3e7e98})       & 1,742 & 3,175 & 5,857 & 9,831 \\
\bottomrule
\end{tabular}
\caption{附录 A.1：强扩展中位 FOM（\texttt{-s 30 -i 100}）}
\end{table}

\subsection{强扩展中位 FOM（\texttt{-s 45 -i 200}，集群 \texttt{cpu-1239}）}

\begin{table}[H]
\centering
\begin{tabular}{lrrrr}
\toprule
\textbf{版本} & \textbf{T=1} & \textbf{T=2} & \textbf{T=4} & \textbf{T=8} \\
\midrule
baseline (\texttt{3e01c40}) & 1,594 & 2,249 & 3,441 & 4,222 \\
v1 (\texttt{3003bb6})       & 1,410 & 1,595 & 1,885 & 1,965 \\
v2 (\texttt{0b4c125})       & 1,247 & 1,475 & 1,790 & 1,914 \\
v3 (\texttt{4d14004})       & 1,493 & 2,379 & 3,904 & 5,454 \\
v9 (\texttt{e3e7e98})       & 1,858 & 3,507 & 5,836 & 7,813 \\
\bottomrule
\end{tabular}
\caption{附录 A.2：强扩展中位 FOM（\texttt{-s 45 -i 200}）}
\end{table}

\subsection{弱扩展中位 FOM（集群 \texttt{cpu-1239}）}

\begin{table}[H]
\centering
\begin{tabular}{lrrrr}
\toprule
\textbf{版本} & \textbf{T=1 (s30)} & \textbf{T=2 (s38)} & \textbf{T=4 (s48)} & \textbf{T=8 (s60)} \\
\midrule
baseline (\texttt{3e01c40}) & 1,479 & 2,220 & 3,536 & 4,102 \\
v1 (\texttt{3003bb6})       & 1,329 & 1,604 & 1,880 & 1,923 \\
v2 (\texttt{0b4c125})       & 1,160 & 1,475 & 1,800 & 1,888 \\
v3 (\texttt{4d14004})       & 1,608 & 2,389 & 4,030 & 5,197 \\
v9 (\texttt{e3e7e98})       & 1,692 & 3,403 & 5,926 & 7,574 \\
\bottomrule
\end{tabular}
\caption{附录 A.3：弱扩展中位 FOM}
\end{table}

\subsection{GPU 性能数据（集群 \texttt{gpu-1239}，\texttt{f2520fa}，CUDA）}

\begin{table}[H]
\centering
\begin{tabular}{lll}
\toprule
\textbf{规模} & \textbf{5 次 FOM} & \textbf{中位数} \\
\midrule
\texttt{-s 30 -i 100} & 11,706 / 14,355 / 14,616 / 14,547 / 14,292 & \textbf{14,355} \\
\texttt{-s 45 -i 200} & 23,046 / 23,460 / 23,433 / 23,327 / 21,277 & \textbf{23,327} \\
\bottomrule
\end{tabular}
\caption{附录 A.4：GPU 性能数据}
\end{table}

\end{document}
